\section{VLM results}
\label{appendix:vlmresults}

\subsection{Annotation}
\label{subsec:vlm_annotation}

We label training data for the oracle with a small web annotation tool built for this purpose.
For each frame, an annotator marks it either as \textit{approach} (is a junction visible, and which of left/straight/right are open) or \textit{exit} (has the robot cleared the junction), or skips it.
Labels are saved as one JSON entry per frame, keyed by filename.


\subsection{Prompting}
\label{subsec:vlm_prompting}

Our zero-shot and few-shot baselines query SmolVLM2 with a fixed, hand-written prompt per question and read off the answer via logit comparison over a closed vocabulary:

\begin{quote}
\textit{Approach:} ``Is there a road junction ahead, and is \{direction\} an open path from it? Answer yes (junction, \{direction\} open), pass (junction, \{direction\} blocked), or no (no junction). One word only.''

\textit{Exit:} ``Has the robot cleared the junction and returned to normal lane-following? Answer yes or no. One word only.''
\end{quote}

For few-shot, we prepend three labeled approach exemplars (one yes, one pass, one no) and two exit exemplars (one yes, one no) as prior turns before the real query.


\subsection{Fine-tuned Classification Head}
\label{subsec:vlm_training}

Instead of prompting, we freeze the SmolVLM2 vision backbone and train a small multi-label MLP head on top of its pooled features to predict the five targets directly: junction, left/straight/right open, and exit-cleared.
We split the image into left/center/right thirds and pool each separately before concatenating, which we found necessary for the model to localize left/right turns rather than only detect a junction.
We compare two input types: the cropped camera image (\textit{Photo}) and the cropped, colorized segmentation mask (\textit{Mask}), and two backbone sizes (SmolVLM2-256M and 500M).
Since our label distribution is imbalanced, we oversample minority-class frames with augmented copies (flip, jitter) and additionally weight the loss per class.
We split by episode rather than by frame to avoid near-duplicate frames leaking between train and validation, and train with Adam and early stopping on validation macro-F1.


\subsection{CNN Baseline}
\label{subsec:vlm_cnn_baseline}

To check whether the VLM backbone is doing useful work, or whether the segmentation mask alone already contains enough signal, we train a small CNN (four conv layers, the same trunk used by our PilotNet models) with the same multi-label head, labels, split, and oversampling as above, directly on the raw segmentation mask.


\subsection{Comparison}
\label{subsec:vlm_comparison}

\begin{table}[!ht]
\centering
\caption{F1 scores for the evaluated approaches.}
\label{tab:classification_results}
\small
\begin{tabular}{lccccc}
\toprule
Metric & Zero-shot & Few-shot & SmolVLM2 & SmolVLM2 & CNN \\
 & Photo & Photo & Mask & Photo & Mask \\
\midrule
Junction        & 0.681 & 0.719 & \textbf{1.000} & 0.970 & 0.836 \\
Left Open       & 0.720 & 0.417 & 0.923 & \textbf{1.000} & 0.720 \\
Straight Open   & 0.222 & 0.000 & \textbf{1.000} & \textbf{1.000} & 0.933 \\
Right Open      & 0.250 & 0.000 & \textbf{1.000} & \textbf{1.000} & 0.609 \\
Exit Cleared    & 0.000 & 0.000 & 0.945 & \textbf{0.963} & 0.612 \\
\midrule
Macro F1        & 0.375 & 0.227 & 0.974 & \textbf{0.987} & 0.742 \\
\bottomrule
\end{tabular}
\end{table}

Zero-shot prompting fails almost completely on exit-cleared (F1 = 0): the question needs a notion of the robot's recent trajectory that a single frame can't provide.
Few-shot exemplars make this worse rather than better, since a handful of in-context images add latency without giving the model a way to localize itself relative to the junction.
Fine-tuning closes this gap: both SmolVLM2 variants clear 0.97 macro F1, with Photo slightly ahead of Mask.
The CNN-Mask baseline, trained on identical labels and input, only reaches 0.742, so the pretrained VLM features are carrying real signal beyond what the mask alone gives a from-scratch CNN.
Despite its slightly lower F1, we deploy the SmolVLM2 Mask head for live inference. The segmentation mask is restricted to our fixed set of classes by construction, while the raw photo can go out-of-distribution whenever objects not seen during training appear in frame.
